\subsection{Sprint 0}

Durante a \textit{sprint} inicial, dediquei minha atenção ao aprendizado das tecnologias de \textit{front-end}. Essa escolha, embora me afastasse das atividades chave desta fase, especificamente a elaboração dos \textit{mockups}, mostrou-se fundamental para adquirir uma visão abrangente sobre a estrutura e os requisitos do projeto. Este período destacou para mim a importância de uma atitude proativa e de uma busca contínua por maneiras de agregar valor à equipe. Entendi que questionar e solicitar esclarecimentos não só enriquece meu entendimento sobre as tarefas em mãos, mas também incentiva uma cultura de colaboração e partilha de saberes entre os membros da equipe.

Este foco único permitiu-me mergulhar nos detalhes de React e TypeScript, com o objetivo de fortalecer significativamente minha capacidade de contribuir para o projeto nas fases futuras. Assim, este esforço não apenas preparou o terreno para uma participação mais efetiva nas próximas \textit{sprints}, mas também sublinhou o valor de uma preparação cuidadosa e de uma participação ativa na construção coletiva do projeto.